\clearpage
\begin{examplebox}
\textbf{\textcolor{CExample}{例 1｜入门：只用照镜子追踪函数值（轴轴）}}
\begin{center}\begin{tikzpicture}[font=\small,>=stealth,line join=round]
\draw[gray!40,rounded corners=3pt] (0,0) rectangle (14.2,8.3);\node at (7.1,7.700000000000001) {例1：轴在1与3，0、2、4的高度都为7};
\begin{scope}[shift={(1.7774193548387096,1.88)},x=0.9677419354838709cm,y=0.45999999999999996cm]\draw[->,gray] (-0.7,0)--(5.5,0) node[right]{$x$};\draw[->,gray] (0,-0.5)--(0,9.5) node[above]{$y$};\draw[blue,thick,solid,domain=-0.3:5,samples=180] plot(\x,{7+cos(deg(pi*(\x-1)/2))});\fill[orange!85!black] (0,7) circle(2.2pt);\fill[orange!85!black] (2,7) circle(2.2pt);\fill[orange!85!black] (4,7) circle(2.2pt);\foreach \a in {1,3}{\draw[violet,dashed] (\a,5.5)--(\a,8.5);\node[above] at (\a,8.5){$x=\a$};}\end{scope}\node[anchor=west,align=left,text width=5.0cm] at (8.4,4.0) {$f=7+\cos\frac{\pi(x-1)}2$\\[9pt]$f(0)=f(2)=f(4)=7$\\[9pt]轴间距2，周期可取4\\[9pt]$2024=506\times4$\\[9pt]$f(2024)=7$};
\node[text width=12.6cm,align=center] at (7.1,.65) {图形与例题初值7相符；两次反射足以求值，无须确定函数解析式。};
\end{tikzpicture}\end{center}



\textbf{题目：}函数 $f(x)$ 在实数范围内有定义，满足
$f(1+t)=f(1-t)$、$f(3+t)=f(3-t)$，并且 $f(0)=7$。求一个周期，并求 $f(2),f(4),f(2024)$。

\textbf{第一步：}两条竖直对称轴分别是 $x=1,x=3$，相距 $2$，所以一个周期是 $T=2\times2=4$。

\textbf{第二步：}关于 $x=1$ 对称，$0$ 与 $2$ 左右对应，故 $f(2)=f(0)=7$。
关于 $x=3$ 对称，$2$ 与 $4$ 对应，故 $f(4)=f(2)=7$。

\textbf{第三步：}$2024=4\times506$，所以 $f(2024)=f(0)=7$。

\textbf{答案：}$\boxed{T=4,\ f(2)=f(4)=f(2024)=7}$。

\medskip

\end{examplebox}
\clearpage
\begin{examplebox}
\textbf{\textcolor{CExample}{例 2｜初中也能算：两个点为什么上下翻两次（点点）}}
\begin{center}\begin{tikzpicture}[font=\small,>=stealth,line join=round]
\draw[gray!40,rounded corners=3pt] (0,0) rectangle (14.2,8.3);\node at (7.1,7.700000000000001) {例2：同高中心翻两次，5变成1，再回到5};
\begin{scope}[shift={(2.4846153846153847,1.9785714285714284)},x=0.7692307692307693cm,y=0.6571428571428571cm]\draw[->,gray] (-1.8,0)--(6,0) node[right]{$x$};\draw[->,gray] (0,-0.5)--(0,6.5) node[above]{$y$};\draw[teal,thick,solid,domain=-1.4:5.6,samples=180] plot(\x,{3+2*sin(deg(pi*\x/2))});\fill[orange!85!black] (1,5) circle(2.2pt);\fill[orange!85!black] (-1,1) circle(2.2pt);\fill[orange!85!black] (5,5) circle(2.2pt);\fill[orange!85!black] (0,3) circle(2.2pt);\fill[orange!85!black] (2,3) circle(2.2pt);\draw[gray,dashed] (-1.6,3)--(5.8,3);\end{scope}\node[anchor=west,align=left,text width=5.0cm] at (8.4,4.0) {$f=3+2\sin\frac{\pi x}2$\\[9pt]中心 $(0,3),(2,3)$\\[9pt]$1\to-1\to5$\\[9pt]高度 $5\to1\to5$\\[9pt]正周期可取4};
\node[text width=12.6cm,align=center] at (7.1,.65) {两个中心共用高度3，纵坐标经过两次关于3的翻转后还原。};
\end{tikzpicture}\end{center}



\textbf{题目：}函数图像关于 $(0,3)$ 与 $(2,3)$ 对称，且 $f(1)=5$。求一个周期和 $f(-1),f(5)$。

\textbf{第一步：}两中心\textbf{高度同为 $3$}，相距 $2$，所以 $T=4$。

\textbf{第二步：}关于 $(0,3)$ 对称，横坐标 $1$ 对应 $-1$；由于两函数值相加为 $6$，有
\[
 f(-1)=6-f(1)=6-5=1.
\]
\textbf{第三步：}关于 $(2,3)$ 对称，$-1$ 对应 $5$，又有
\[
 f(5)=6-f(-1)=6-1=5.
\]
所以 $f(5)=f(1)$，正好右移 $4$ 后恢复。\textbf{答案：}$T=4,\ f(-1)=1,\ f(5)=5$。

\medskip

\end{examplebox}
\clearpage
\begin{examplebox}
\textbf{\textcolor{CExample}{例 3｜一眼理解“四倍”：先变低，再恢复（轴点）}}
\begin{center}\begin{tikzpicture}[font=\small,>=stealth,line join=round]
\draw[gray!40,rounded corners=3pt] (0,0) rectangle (14.2,8.3);\node at (7.1,7.700000000000001) {例3：轴在0，中心 $(1,3)$；走2高度尚未还原};
\begin{scope}[shift={(2.174626865671642,1.9785714285714284)},x=0.8955223880597015cm,y=0.6571428571428571cm]\draw[->,gray] (-1.2,0)--(5.5,0) node[right]{$x$};\draw[->,gray] (0,-0.5)--(0,6.5) node[above]{$y$};\draw[blue,thick,solid,domain=-0.8:5,samples=180] plot(\x,{3+2*cos(deg(pi*\x/2))});\fill[orange!85!black] (0,5) circle(2.2pt);\fill[orange!85!black] (2,1) circle(2.2pt);\fill[orange!85!black] (4,5) circle(2.2pt);\fill[orange!85!black] (1,3) circle(2.2pt);\draw[violet,dashed] (0,-.3)--(0,6);\draw[gray,dashed] (-1,3)--(5.2,3);\end{scope}\node[anchor=west,align=left,text width=5.0cm] at (8.4,4.0) {$f=3+2\cos\frac{\pi x}2$\\[9pt]轴 $x=0$，中心 $(1,3)$\\[9pt]$f(0)=5,\ f(2)=1$\\[9pt]$f(4)=5$\\[9pt]$d=1,\ T=4$};
\node[text width=12.6cm,align=center] at (7.1,.65) {相对高度3一高一低；移动2只翻转，移动4才保证原样重复。};
\end{tikzpicture}\end{center}



\textbf{题目：}$f(t)=f(-t)$，$f(1+t)+f(1-t)=6$，并给出 $f(0)=5$。求一个周期与 $f(2),f(4),f(100)$。

第一条表示轴为 $x=0$，第二条表示中心是 $(1,3)$。距离 $d=1$，所以一个周期是 $T=4$。

由本结论的关键一步：$f(x+2)=6-f(x)$，因此
\[
 f(2)=6-f(0)=1,\qquad f(4)=6-f(2)=5.
\]
又 $100=4\times25$，所以 $f(100)=f(0)=5$。

\textbf{关键发现：}$f(0)=5$、$f(2)=1$、$f(4)=5$。\textbf{隔 $2$ 只变成围绕 $3$ 的另一侧，隔 $4$ 才真正相同。}

\medskip

\end{examplebox}
\clearpage
\begin{examplebox}
\textbf{\textcolor{CExample}{例 4｜高考常见：不求表达式，直接算很大的函数值}}
\begin{center}\begin{tikzpicture}[font=\small,>=stealth,line join=round]
\draw[gray!40,rounded corners=3pt] (0,0) rectangle (14.2,8.3);\node at (7.1,7.700000000000001) {例4：同高中心在1与3，周期4把2025送回1};
\begin{scope}[shift={(1.7774193548387096,2.1727272727272724)},x=0.9677419354838709cm,y=0.5227272727272726cm]\draw[->,gray] (-0.7,0)--(5.5,0) node[right]{$x$};\draw[->,gray] (0,-1)--(0,7.8) node[above]{$y$};\draw[teal,thick,solid,domain=-0.3:5,samples=180] plot(\x,{3-3*cos(deg(pi*\x/2))});\fill[orange!85!black] (1,3) circle(2.2pt);\fill[orange!85!black] (3,3) circle(2.2pt);\fill[orange!85!black] (0,0) circle(2.2pt);\fill[orange!85!black] (4,0) circle(2.2pt);\draw[gray,dashed] (-.5,3)--(5.2,3);\end{scope}\node[anchor=west,align=left,text width=5.0cm] at (8.4,4.0) {$f=3-3\cos\frac{\pi x}2$\\[9pt]中心 $(1,3),(3,3)$\\[9pt]$f(0)=f(4)=0$\\[9pt]$f(1)=3$\\[9pt]$f(2025)=f(1)=3$};
\node[text width=12.6cm,align=center] at (7.1,.65) {中心所在的横坐标代入后函数值为3；同高中心相距2，给出周期4。};
\end{tikzpicture}\end{center}



\textbf{题目：}设 $f(1+t)+f(1-t)=6$、$f(3+t)+f(3-t)=6$，且 $f(0)=0$。求一个周期以及 $f(2025)$、$f(4)$。

两个对称中心分别为 $(1,3),(3,3)$，同高度且相距 $2$，因此 $T=4$。

$2025=4\times506+1$，因此 $f(2025)=f(1)$。在第一条等式中令 $t=0$，有 $2f(1)=6$，故 $f(1)=3$，于是 $f(2025)=3$。

又 $f(4)=f(0)=0$。\textbf{答案：}$\boxed{T=4,\ f(2025)=3,\ f(4)=0}$。

\medskip

\end{examplebox}
\clearpage
\begin{examplebox}
\textbf{\textcolor{CExample}{例 5｜稍难：把轴点关系推广到任意高度}}
\begin{center}\begin{tikzpicture}[font=\small,>=stealth,line join=round]
\draw[gray!40,rounded corners=3pt] (0,0) rectangle (14.2,8.3);\node at (7.1,7.700000000000001) {例5：相隔8先绕高度4翻转，相隔16才恢复};
\begin{scope}[shift={(1.330769230769231,3.419230769230769)},x=0.23076923076923078cm,y=0.11794871794871795cm]\draw[->,gray] (-1,0)--(25,0) node[right]{$x$};\draw[->,gray] (0,-15)--(0,24) node[above]{$y$};\draw[blue,thick,solid,domain=0:24,samples=180] plot(\x,{4+6*cos(deg(pi*(\x-2)/8))/cos(deg(3*pi/8))});\fill[orange!85!black] (5,10) circle(2.2pt);\fill[orange!85!black] (13,-2) circle(2.2pt);\fill[orange!85!black] (21,10) circle(2.2pt);\fill[orange!85!black] (6,4) circle(2.2pt);\draw[gray,dashed] (0,4)--(24,4);\draw[violet,dashed] (2,-13)--(2,21);\foreach \v in {5,13,21}{\node[below=4pt] at (\v,-13){$\v$};}\end{scope}\node[anchor=west,align=left,text width=5.0cm] at (8.4,4.0) {轴 $x=2$，中心 $(6,4)$\\[9pt]图取 $f=4+A\cos\frac{\pi(x-2)}8$\\[9pt]$A=6/\cos(3\pi/8)$\\[9pt]$f(5)=10,\ f(13)=-2$\\[9pt]$f(21)=10,\ T=16$};
\node[text width=12.6cm,align=center] at (7.1,.65) {振幅按初值精确确定；同一条曲线上，10与负2关于高度4成对。};
\end{tikzpicture}\end{center}



\textbf{题目：}已知 $f(2+t)=f(2-t)$，$f(6+t)+f(6-t)=8$。若 $f(5)=10$，求一个周期以及 $f(13),f(21)$。

轴在 $2$，中心为 $(6,4)$，相距 $4$，所以一个周期 $T=16$。

但先不急着用 $16$；两处对称合起来，\textbf{每隔 $8$，函数值会变为 $8-f(x)$}：
\[
 f(x+8)=8-f(x).
\]
于是 $f(13)=8-f(5)= -2$；$f(21)=8-f(13)=10$。

\textbf{答案：}$\boxed{T=16,\ f(13)=-2,\ f(21)=10}$。

\medskip

\end{examplebox}
\clearpage
\begin{examplebox}
\textbf{\textcolor{CExample}{例 6｜很容易掉坑：“一个周期”不是“最小正周期”}}
\begin{center}\begin{tikzpicture}[font=\small,>=stealth,line join=round]
\draw[gray!40,rounded corners=3pt] (0,0) rectangle (14.2,8.3);\node at (7.1,7.700000000000001) {例6：两轴相距2，只保证4是周期，未必最小};
\begin{scope}[shift={(1.58,3.814705882352941)},x=1.2cm,y=1.352941176470588cm]\draw[->,gray] (-0.4,0)--(4.6,0) node[right]{$x$};\draw[->,gray] (0,-1.6)--(0,1.8) node[above]{$y$};\draw[blue,thick,solid,domain=-0.1:4.2,samples=180] plot(\x,{cos(deg(3*pi*\x/2))});\fill[orange!85!black] (0,1) circle(2.2pt);\fill[orange!85!black] (1.3333333333333333,1) circle(2.2pt);\fill[orange!85!black] (2.6666666666666665,1) circle(2.2pt);\fill[orange!85!black] (4,1) circle(2.2pt);\end{scope}\node[anchor=west,align=left,text width=5.0cm] at (8.4,4.0) {$f=\cos\frac{3\pi x}2$\\[9pt]轴 $x=0,2$\\[9pt]公式给出周期4\\[9pt]最小正周期 $\frac43$\\[9pt]4含三个完整重复};
\node[text width=12.6cm,align=center] at (7.1,.65) {具体曲线在公式给出的距离内已经重复三次；最小周期需要额外判断。};
\end{tikzpicture}\end{center}



\textbf{题目：}设 $f(x)=\cos\!\left(\dfrac{3\pi x}{2}\right)$，指出它的两条对称轴，并比较公式周期与最小正周期。

由于余弦在 $x=0$ 和 $x=2$ 处左右对称，轴轴公式得出\textbf{一个周期} $T=2\times2=4$。

但余弦内部每增加 $2\pi$，函数值就重复。因此
\[
\frac{3\pi}{2}T_{\min}=2\pi
\quad\Longrightarrow\quad T_{\min}=\frac43.
\]

$4=3\times\frac43$，所以 $4$ 当然也是周期。\textbf{题目若只求“一个周期”，答 $4$ 正确；若求“最小正周期”，就不能只套双对称公式。}

\medskip

\end{examplebox}
\clearpage
\begin{examplebox}
\textbf{\textcolor{CExample}{例 7｜反例：两个中心不在同一高度，为什么不一定有周期？}}
\begin{center}\begin{tikzpicture}[font=\small,>=stealth,line join=round]
\draw[gray!40,rounded corners=3pt] (0,0) rectangle (14.2,8.3);\node at (7.1,7.700000000000001) {例7：两个中心一高一低，直线只会斜着搬移};
\begin{scope}[shift={(2.225,2.540322580645161)},x=0.75cm,y=0.49462365591397844cm]\draw[->,gray] (-1.5,0)--(6.5,0) node[right]{$x$};\draw[->,gray] (0,-1.8)--(0,7.5) node[above]{$y$};\draw[blue,thick,solid,domain=-1:6,samples=180] plot(\x,{\x});\fill[orange!85!black] (0,0) circle(2.2pt);\fill[orange!85!black] (2,2) circle(2.2pt);\fill[orange!85!black] (1,1) circle(2.2pt);\fill[orange!85!black] (5,5) circle(2.2pt);\draw[orange!85!black,dashed] (1,1)--(5,5);\end{scope}\node[anchor=west,align=left,text width=5.0cm] at (8.4,4.0) {$f=x$\\[9pt]中心 $(0,0),(2,2)$\\[9pt]$x\to x+4$ 时\\[9pt]$f(x+4)=f(x)+4$\\[9pt]不保持原高度};
\node[text width=12.6cm,align=center] at (7.1,.65) {两次中心对称会同时改变横纵坐标；同高度条件不可省略。};
\end{tikzpicture}\end{center}



令 $f(x)=x$，这条斜直线关于原点 $(0,0)$ 对称，也关于 $(2,2)$ 对称，因为
\[
 f(t)+f(-t)=0,\qquad f(2+t)+f(2-t)=4.
\]

虽然这两个中心横坐标相差 $2$，但高度一处为 $0$、另一处为 $2$；$f(x+4)=x+4\ne f(x)$。

\textbf{因此“点点周期 $2d$”中的“两个中心同高度”不可省略。}
\end{examplebox}
